\documentclass[10pt,a4paper]{article}

\usepackage[margin=1.35cm]{geometry}
\usepackage{fontspec}
\usepackage{xcolor}
\usepackage{enumitem}
\usepackage[hidelinks]{hyperref}
\usepackage{titlesec}
\usepackage{tabularx}
\usepackage{array}
\usepackage{ragged2e}

\setmainfont{TeX Gyre Heros}
\setsansfont{TeX Gyre Heros}
\renewcommand{\familydefault}{\sfdefault}
\definecolor{accent}{HTML}{17324D}
\definecolor{muted}{HTML}{5F6B76}
\definecolor{rulegray}{HTML}{D8DEE5}
\definecolor{bodytext}{HTML}{1F2933}

\pagestyle{empty}
\setlength{\parindent}{0pt}
\setlength{\parskip}{0.25em}
\setlist[itemize]{leftmargin=1.4em,itemsep=0.25em,topsep=0.2em}
\titleformat{\section}
  {\large\bfseries\color{accent}}
  {}
  {0pt}
  {}
  [\vspace{-0.25em}\color{rulegray}\titlerule]
\titlespacing*{\section}{0pt}{0.8em}{0.45em}

\newcommand{\tagline}[1]{\textcolor{muted}{\large #1}}
\newcommand{\role}[1]{\textbf{\color{bodytext}#1}}
\newcommand{\venue}[1]{\textcolor{muted}{#1}}
\newcommand{\pubitem}[3]{%
  \textbf{#1}\hfill\venue{#3}\\
  #2\par\vspace{0.35em}
}

\begin{document}

{\Huge\bfseries\color{accent} Yuxiang Lin}\\[0.3em]
\tagline{Ph.D. Candidate in Computational Biology and Bioinformatics}\\
\tagline{Xiamen University \quad | \quad Expected 2027}\\[0.55em]
{\large Seeking 2027 postdoc / research fellow / research scientist interview opportunities}\\
{\large Preferred regions: \textbf{Singapore} and \textbf{North America}}\\[0.45em]
\textcolor{muted}{Public profile: \href{https://lenislin.github.io/}{https://lenislin.github.io/}}

\vspace{0.8em}
\color{rulegray}\hrule
\vspace{0.75em}
\color{bodytext}

\section*{Research Fit}
\begin{tabularx}{\textwidth}{>{\bfseries\color{accent}}p{0.24\textwidth} X}
Spatial transcriptomics &
Foundation modeling and predictive learning for spatial molecular profiles, with emphasis on tissue context and image-gene alignment.\\[0.45em]
Tumor microenvironment &
Computational discovery of clinically relevant niches, recurrence-risk signatures, and translational biomarkers.\\[0.45em]
Biomedical AI agents &
Tool-augmented biomedical research workflows involving literature-aware reasoning, multimodal evidence integration, and automated analysis.\\[0.45em]
Multi-omics medicine &
Integration of single-cell, spatial, transcriptomic, proteomic, and imaging data for disease mechanism discovery.
\end{tabularx}

\section*{Selected Work}
\pubitem{ST-Align}
{Multi-scale image-gene foundation modeling for spatial transcriptomics via spot-niche alignment.}
{ICLR Workshop, 2026}

\pubitem{TiRank}
{Prioritizes phenotypic niches in the tumor microenvironment for clinical biomarker discovery.}
{Genome Medicine, 2026}

\pubitem{HyperST}
{Hierarchical hyperbolic learning for spatial transcriptomics prediction.}
{CVPR Spotlight, 2026}

\pubitem{Single-cell recurrence modeling}
{Prioritization of prognostic-associated subpopulations and individualized recurrence-risk signatures from colorectal cancer single-cell transcriptomes.}
{Briefings in Bioinformatics, 2023}

\section*{Academic Service}
Reviewer experience includes ICLR, CVPR, IEEE BIBM, Nature Communications, IEEE Transactions on Medical Imaging, Bioinformatics, and Briefings in Bioinformatics.

\section*{What I Am Looking For}
I am open to early discussions with labs working on computational biology, spatial omics, biomedical foundation models, clinical translation, or agentic AI systems for biomedical research. Target start time: 2027.

\vfill
{\small\textcolor{muted}{Public version for opportunity discovery. Full CV and references can be shared after mutual fit is established.}}

\end{document}
